\subsection*{پرسش ۱: یادگیری مدل انرژی با \lr{Score Matching} (۲۰ نمره)}

فرض کنید از چگالی ناشناختهٔ $p_{\mathrm{data}}(x)>0$ روی $\mathbb{R}^{d}$ نمونه داریم و می‌خواهیم مدل زیر را آموزش دهیم:
\begin{flalign*}
p_{\theta}(x)
&=\frac{\exp(-E_{\theta}(x))}{Z(\theta)},
\qquad
Z(\theta)=\int_{\mathbb{R}^{d}}\exp(-E_{\theta}(u))\,du,
&&\\
s_{\theta}(x)
&=\nabla_x\log p_{\theta}(x).
&&
\end{flalign*}
محاسبهٔ $Z(\theta)$ ممکن است دشوار باشد. هدف اولیهٔ \lr{Score Matching} را به صورت زیر تعریف می‌کنیم:
\begin{flalign*}
\mathcal{L}(\theta)
&=\frac{1}{2}\,
\mathbb{E}_{x\sim p_{\mathrm{data}}}
\left[
\left\|s_{\theta}(x)-\nabla_x\log p_{\mathrm{data}}(x)\right\|_2^2
\right].
&&
\end{flalign*}
فرض کنید توابع به اندازهٔ کافی هموارند، تمام امیدهای مورد استفاده متناهی‌اند و جمله‌های مرزی در انتگرال‌گیری جزءبه‌جزء صفر می‌شوند.

\begin{enumerate}
\item[الف)] (۴ نمره)
$s_{\theta}(x)$ را برحسب $E_{\theta}(x)$ به دست آورید.
چرا وابستگی $Z$ به $\theta$ در این محاسبه مشکل‌ساز نیست؟
آیا هنگام محاسبهٔ $\nabla_{\theta}\log p_{\theta}(x)$ نیز می‌توان جملهٔ مربوط به $Z(\theta)$ را حذف کرد؟

\item[ب)] (۸ نمره)
هدف اولیه هنوز به $\nabla_x\log p_{\mathrm{data}}(x)$ وابسته است که در اختیار نداریم.
ابتدا مربع را بسط دهید؛ سپس با انتگرال‌گیری جزءبه‌جزء نشان دهید که می‌توان هزینه‌ای با همان کمینه‌سازها به دست آورد که فقط شامل
$s_{\theta}(x)$ و $\operatorname{div}s_{\theta}(x)$ باشد.
این هزینه را برحسب $E_{\theta}$ نیز بنویسید و نسخهٔ تجربی آن را برای نمونه‌های
$x^{(1)},\ldots,x^{(n)}$ مشخص کنید.

راهنمایی: جملهٔ مختلط را به مجموع عبارت‌های
$\int s_{\theta,i}(x)\,\partial_{x_i}p_{\mathrm{data}}(x)\,dx$
تبدیل کنید. از تعاریف زیر استفاده کنید:
\begin{flalign*}
\operatorname{div}s_{\theta}(x)
&=\sum_{i=1}^{d}\frac{\partial s_{\theta,i}(x)}{\partial x_i},
\qquad
\Delta E_{\theta}(x)
=\sum_{i=1}^{d}\frac{\partial^2 E_{\theta}(x)}{\partial x_i^2}.
&&
\end{flalign*}

\item[ج)] (۸ نمره)
اکنون $d=1$ و مدل انرژی زیر را در نظر بگیرید:
\begin{flalign*}
E_{\theta}(x)
&=\frac{\theta}{2}x^2+\frac{\lambda}{4}x^4,
\qquad \theta\in\mathbb{R},\quad \lambda>0,
&&\\
m_k
&=\frac{1}{n}\sum_{j=1}^{n}\bigl(x^{(j)}\bigr)^k,
\qquad k\in\{2,4,6\},\quad m_2>0.
&&
\end{flalign*}
$\lambda$ معلوم است و فقط $\theta$ آموزش داده می‌شود.
با استفاده از بخش قبل، هزینهٔ تجربی و کمینه‌ساز یکتای
$\widehat{\theta}(\lambda)$ را به دست آورید.

سپس داده‌ها را ثابت نگه دارید و فرض کنید برای هر مقدار $\lambda$، مدل دوباره برازش می‌شود.
آستانهٔ تغییر چگالی برازش‌شده از تک‌قله‌ای به دوقله‌ای و مکان قله‌ها را مشخص کنید.
آیا منفی‌شدن $\widehat{\theta}$ باعث غیرقابل‌نرمال‌سازی‌شدن مدل می‌شود؟
\end{enumerate}

\subsection*{پاسخ‌نامهٔ پرسش ۱}

\subsubsection*{الف)}
\begin{flalign*}
s_{\theta}(x)
&=\nabla_x\bigl[-E_{\theta}(x)-\log Z(\theta)\bigr]
=-\nabla_x E_{\theta}(x),
&&\\
\nabla_x\log Z(\theta)
&=0,
&&\\
\nabla_{\theta}\log p_{\theta}(x)
&=-\nabla_{\theta}E_{\theta}(x)
-\nabla_{\theta}\log Z(\theta).
&&
\end{flalign*}
$Z$ نسبت به $x$ ثابت است؛ اما مشتق آن نسبت به $\theta$ عموماً صفر نیست.

\subsubsection*{ب)}
با انتگرال‌گیری جزءبه‌جزء:
\begin{flalign*}
\mathbb{E}_{p_{\mathrm{data}}}
\left[s_{\theta}^{\top}\nabla_x\log p_{\mathrm{data}}\right]
&=\sum_{i=1}^{d}
\int_{\mathbb{R}^{d}}
s_{\theta,i}(x)\,\partial_{x_i}p_{\mathrm{data}}(x)\,dx
&&\\
&=-\sum_{i=1}^{d}
\int_{\mathbb{R}^{d}}
p_{\mathrm{data}}(x)\,\partial_{x_i}s_{\theta,i}(x)\,dx
&&\\
&=-\mathbb{E}_{p_{\mathrm{data}}}
\left[\operatorname{div}s_{\theta}(x)\right].
&&
\end{flalign*}
بنابراین:
\begin{flalign*}
\mathcal{L}(\theta)
&=
\mathbb{E}_{p_{\mathrm{data}}}
\left[
\frac{1}{2}\|s_{\theta}(x)\|_2^2
+\operatorname{div}s_{\theta}(x)
\right]+C,
&&\\
C
&=\frac{1}{2}\mathbb{E}_{p_{\mathrm{data}}}
\left[\|\nabla_x\log p_{\mathrm{data}}(x)\|_2^2\right],
\qquad \nabla_{\theta}C=0,
&&\\
\mathcal{J}(\theta)
&=
\mathbb{E}_{p_{\mathrm{data}}}
\left[
\frac{1}{2}\|\nabla_x E_{\theta}(x)\|_2^2
-\Delta E_{\theta}(x)
\right],
&&\\
\widehat{\mathcal{J}}(\theta)
&=\frac{1}{n}\sum_{j=1}^{n}
\left[
\frac{1}{2}\|\nabla_x E_{\theta}(x^{(j)})\|_2^2
-\Delta E_{\theta}(x^{(j)})
\right].
&&
\end{flalign*}
$\mathcal{J}$ با $\mathcal{L}$ کمینه‌سازهای یکسان دارد؛
$\widehat{\mathcal{J}}$ تقریب نمونه‌ای آن است.

\subsubsection*{ج)}
\begin{flalign*}
s_{\theta}(x)
&=-\theta x-\lambda x^3,
\qquad
s_{\theta}'(x)=-\theta-3\lambda x^2,
&&\\
\widehat{\mathcal{J}}(\theta)
&=\frac{\theta^2}{2}m_2
+\theta\lambda m_4
+\frac{\lambda^2}{2}m_6
-\theta-3\lambda m_2,
&&\\
\frac{d\widehat{\mathcal{J}}}{d\theta}
&=\theta m_2+\lambda m_4-1,
\qquad
\frac{d^2\widehat{\mathcal{J}}}{d\theta^2}=m_2>0,
&&\\
\boxed{\widehat{\theta}(\lambda)}
&=\boxed{\frac{1-\lambda m_4}{m_2}}.
&&
\end{flalign*}

قله‌های چگالی، کمینه‌های انرژی‌اند:
\begin{flalign*}
E_{\widehat{\theta}}'(x)
&=x\bigl(\widehat{\theta}+\lambda x^2\bigr),
&&\\
\lambda_{\mathrm{c}}
&=\frac{1}{m_4},
&&\\
\operatorname*{arg\,max}_{x}p_{\widehat{\theta}}(x)
&=
\begin{cases}
\{0\}, & 0<\lambda\leq\lambda_{\mathrm{c}},\\[2mm]
\displaystyle
\left\{
-\sqrt{\frac{\lambda m_4-1}{\lambda m_2}},
\sqrt{\frac{\lambda m_4-1}{\lambda m_2}}
\right\},
& \lambda>\lambda_{\mathrm{c}}.
\end{cases}
&&\\
\frac{d\widehat{\theta}}{d\lambda}
&=-\frac{m_4}{m_2}<0.
&&
\end{flalign*}
در آستانه، انرژی برابر $\lambda x^4/4$ است و همچنان فقط یک کمینه در صفر دارد.
منفی‌شدن $\widehat{\theta}$ مانع نرمال‌سازی نیست:
\begin{flalign*}
\lim_{|x|\to\infty}\frac{E_{\theta}(x)}{x^4}
&=\frac{\lambda}{4}>0
\quad\Longrightarrow\quad
0<Z(\theta)<\infty
\qquad(\theta\in\mathbb{R}).
&&
\end{flalign*}

% منبع رابطهٔ پایهٔ پرسش ۲: [Vincent, 2011](https://gwern.net/doc/ai/nn/diffusion/2011-vincent.pdf).
% اقتباس آموزشی: هم‌ارزی DSM و Score Matching به کمک امید شرطی صورت‌بندی شده؛ رابطه با پیش‌بینی نویز و تحلیل یک توزیع دو‌نقطه‌ای برای بررسی اثر شدت نویز به آن افزوده شده است.

\subsection*{پرسش ۲: \lr{Denoising Score Matching} چه چیزی را یاد می‌گیرد؟ (۲۰ نمره)}

برای آموزش شبکهٔ امتیاز، به دادهٔ پاک $X\in\mathbb{R}^{d}$ نویز اضافه می‌کنیم:
\begin{flalign*}
Y&=X+\sigma\varepsilon,
\qquad
\varepsilon\sim\mathcal{N}(0,I_d),
\qquad \varepsilon\perp X,\quad \sigma>0,
&&\\
q_{\sigma}(y\mid x)
&=\frac{1}{(2\pi\sigma^2)^{d/2}}
\exp\left(-\frac{\|y-x\|_2^2}{2\sigma^2}\right),
&&\\
q_{\sigma}(y)
&=\mathbb{E}_{X}\left[q_{\sigma}(y\mid X)\right],
\qquad
s_{\sigma}(y)=\nabla_y\log q_{\sigma}(y).
&&
\end{flalign*}
در این پرسش $\sigma$ ثابت است. فرض کنید مشتق‌گیری زیر امید مجاز است، امیدهای مورد استفاده متناهی‌اند و شبکه ظرفیت کافی برای نمایش تابع بهینه دارد.

\begin{enumerate}
\item[الف)] (۶ نمره)
ابتدا
$h_{\sigma}(y,x)=\nabla_y\log q_{\sigma}(y\mid x)$
را محاسبه کنید. سپس نشان دهید:
\begin{flalign*}
s_{\sigma}(y)
&=\mathbb{E}\left[h_{\sigma}(Y,X)\mid Y=y\right]
=\frac{\mathbb{E}[X\mid Y=y]-y}{\sigma^2}.
&&
\end{flalign*}
راهنمایی: از
$\nabla_y q_{\sigma}(y)
=\mathbb{E}_X[\nabla_y q_{\sigma}(y\mid X)]$
شروع کنید و وزن‌های حاصل را به صورت توزیع شرطی $X\mid Y=y$ تفسیر کنید.

\item[ب)] (۸ نمره)
شبکهٔ $s_{\theta}(y,\sigma)$ را با هزینهٔ زیر آموزش می‌دهیم:
\begin{flalign*}
\mathcal{L}_{D}(\theta)
&=\frac{1}{2}\mathbb{E}_{X,Y}
\left[
\|s_{\theta}(Y,\sigma)-h_{\sigma}(Y,X)\|_2^2
\right].
&&
\end{flalign*}
با استفاده از بخش قبل نشان دهید:
\begin{flalign*}
\mathcal{L}_{D}(\theta)
&=\frac{1}{2}\mathbb{E}_{Y}
\left[
\|s_{\theta}(Y,\sigma)-s_{\sigma}(Y)\|_2^2
\right]+C_{\sigma},
&&
\end{flalign*}
که $C_{\sigma}$ مستقل از $\theta$ است. در نتیجه، تابع بهینهٔ شبکه چیست؟

اکنون پارامترسازی
$s_{\theta}(y,\sigma)=-\varepsilon_{\theta}(y,\sigma)/\sigma$
را در نظر بگیرید.
هزینه را برحسب خطای پیش‌بینی نویز بنویسید و
$\varepsilon^{*}(y,\sigma)$ را مشخص کنید.
آیا برای هر زوج آموزشی $(x,y)$ می‌توان
$s_{\sigma}(y)=-\varepsilon/\sigma$
را یک تساوی دقیق دانست؟

راهنمایی: اختلاف شبکه و هدف را به دو جمله، با واسطهٔ
$s_{\sigma}(Y)$، بشکنید و امید شرطی جملهٔ مختلط را بررسی کنید.

\item[ج)] (۶ نمره)
برای تحلیل اثر شدت نویز، فرض کنید $d=1$ و دادهٔ پاک فقط دو مقدار دارد:
\begin{flalign*}
\mathbb{P}(X=a)
&=\mathbb{P}(X=-a)=\frac{1}{2},
\qquad a>0,
&&\\
q_{\sigma}(y)
&=\frac{1}{2\sqrt{2\pi}\sigma}
\left[
\exp\left(-\frac{(y-a)^2}{2\sigma^2}\right)
+
\exp\left(-\frac{(y+a)^2}{2\sigma^2}\right)
\right].
&&
\end{flalign*}
$s_{\sigma}(y)$ و $s_{\sigma}'(0)$ را به دست آورید.
برای دو حالت $\sigma<a$ و $\sigma>a$، مشخص کنید صفر کمینهٔ محلی چگالی است یا بیشینهٔ محلی، و میدان امتیاز در نزدیکی صفر به کدام جهت اشاره می‌کند.
آیا صفرشدن امتیاز الزاماً به معنای قرارداشتن در قلهٔ چگالی است؟

راهنمایی: از
$e^{u}+e^{-u}=2\cosh u$
و
$\frac{d}{du}\log\cosh u=\tanh u$
استفاده کنید.
\end{enumerate}

\subsection*{پاسخ‌نامهٔ پرسش ۲}

\subsubsection*{الف)}
\begin{flalign*}
h_{\sigma}(y,x)
&=-\frac{y-x}{\sigma^2},
&&\\
s_{\sigma}(y)
&=\frac{\mathbb{E}_{X}
\left[q_{\sigma}(y\mid X)h_{\sigma}(y,X)\right]}
{q_{\sigma}(y)}
&&\\
&=\mathbb{E}\left[h_{\sigma}(y,X)\mid Y=y\right]
&&\\
&=\boxed{\frac{\mathbb{E}[X\mid Y=y]-y}{\sigma^2}}.
&&
\end{flalign*}

\subsubsection*{ب)}
\begin{flalign*}
\mathbb{E}\left[h_{\sigma}(Y,X)-s_{\sigma}(Y)\mid Y\right]
&=0,
&&\\
\mathbb{E}\left[
\bigl(s_{\theta}(Y,\sigma)-s_{\sigma}(Y)\bigr)^{\top}
\bigl(s_{\sigma}(Y)-h_{\sigma}(Y,X)\bigr)
\right]
&=0,
&&\\
\mathcal{L}_{D}(\theta)
&=\frac{1}{2}\mathbb{E}
\left[\|s_{\theta}(Y,\sigma)-s_{\sigma}(Y)\|_2^2\right]
+C_{\sigma},
&&\\
C_{\sigma}
&=\frac{1}{2}\mathbb{E}
\left[\|h_{\sigma}(Y,X)-s_{\sigma}(Y)\|_2^2\right],
&&\\
s^{*}(y,\sigma)
&=s_{\sigma}(y)
=\nabla_y\log q_{\sigma}(y).
&&
\end{flalign*}
این برابری برای تابع بهینه، تقریباً همه‌جا نسبت به $q_{\sigma}$ برقرار است.

با پارامترسازی پیش‌بینی نویز:
\begin{flalign*}
h_{\sigma}(Y,X)
&=-\frac{\varepsilon}{\sigma},
&&\\
\mathcal{L}_{D}(\theta)
&=\frac{1}{2\sigma^2}
\mathbb{E}_{X,\varepsilon}
\left[
\|\varepsilon_{\theta}(X+\sigma\varepsilon,\sigma)-\varepsilon\|_2^2
\right],
&&\\
\varepsilon^{*}(y,\sigma)
&=\mathbb{E}[\varepsilon\mid Y=y]
=-\sigma s_{\sigma}(y).
&&
\end{flalign*}
$-\varepsilon/\sigma$ هدف یک زوج آموزشی است؛
$s_{\sigma}(y)$ میانگین شرطی این هدف است.
این دو عموماً برای یک زوج منفرد برابر نیستند.

\subsubsection*{ج)}
\begin{flalign*}
q_{\sigma}(y)
&=\frac{1}{\sqrt{2\pi}\sigma}
\exp\left(-\frac{y^2+a^2}{2\sigma^2}\right)
\cosh\left(\frac{ay}{\sigma^2}\right),
&&\\
\boxed{s_{\sigma}(y)}
&=\boxed{
-\frac{y}{\sigma^2}
+\frac{a}{\sigma^2}\tanh\left(\frac{ay}{\sigma^2}\right)
},
&&\\
s_{\sigma}'(y)
&=-\frac{1}{\sigma^2}
+\frac{a^2}{\sigma^4}
\left[
1-\tanh^2\left(\frac{ay}{\sigma^2}\right)
\right],
&&\\
s_{\sigma}(0)
&=0,
\qquad
s_{\sigma}'(0)=\frac{a^2-\sigma^2}{\sigma^4},
&&\\
\frac{q_{\sigma}''(0)}{q_{\sigma}(0)}
&=s_{\sigma}'(0)+s_{\sigma}(0)^2
=\frac{a^2-\sigma^2}{\sigma^4}.
&&
\end{flalign*}

برای $\sigma<a$، صفر کمینهٔ محلی چگالی است و میدان امتیاز در نزدیکی آن از صفر دور می‌شود.
برای $\sigma>a$، صفر بیشینهٔ محلی است و میدان به سمت صفر اشاره می‌کند.
پس افزایش نویز می‌تواند مرکزِ میان دو خوشه را به یک قله تبدیل کند؛
$s_{\sigma}(0)=0$ به‌تنهایی میان قله و دره تمایز نمی‌گذارد.